This study has a few limitations.

First, even though DistilBERT improves the test performance overall, class boundaries stay difficult.
The weakest one of them is \textit{Business} with an F1 score of 0.9191. It often gets confused with \textit{Sci/Tech} and vice versa. From the confusion it is visible that \textit{Business} being misclassified as \textit{Sci/Tech} occurred 131 times, while \textit{Sci/Tech} being misclassified as \textit{Business} occurred 83 times. The same pattern is visible for LSTM, where it is even stronger (145 and 141 misclassifications respectively). This suggests that this is not just an issue of model capacity, but maybe also because of overlapping language between finance and tech in the dataset causing ambiguity.

Second, both LSTM and DistilBERT are sensitive to shorter inputs. Macro-F1 drops to 0.8815 and 0.9302 for LSTM and DistilBERT respectively in the shortest length bucket. Both models peak at the mid length bucket (0.9287 and 0.9571 respectively). This is an indication that both models rely heavily on context and are not as reliable when article text is relatively shorter.

Third, the keyword masking probe indicates that both models rely on lexical cues, especially LSTM. KKeyword masking reduced the macro-F1 score more than frequency-matched random masking did. (For DistilBERT 0.0090 vs 0.0009 and LSTM 0.0142 vs 0.0010).

Fourth: robustness analyses have some limitations: the length bucket analysis is based on word count rather than token count, which may not perfectly capture the model's handling of input length, especially for models with different tokenization strategies.

Finally, we only evaluate on AG News, so our conclusions may not transfer to datasets with different writing styles or label structures.

Some future improvements could include evaluating on additional datasets to test the generality of our findings, using token count for length-based analyses, exploring more sophisticated methods for identifying important keywords for the masking probe, and experimenting with different pooling strategies for the LSTM. Additionally, comparing a wider range of encoder only transformer architectures (e.g. BERT LARGE and RoBERTa) could help in finding the best model for this task.